\section{Analiza niepewności}
\label{sec:niepewnosc}
Dla $n$ niezależnych powtórzeń tej samej wielkości, przy założeniu
stabilności warunków, średnią i jej niepewność typu A oblicza się jako
\begin{equation}
  \overline{L}_p = \frac{1}{n}\sum_{i=1}^{n} L_{p,i},
  \qquad
  u_A(\overline{L}_p) =
  \sqrt{\frac{\sum_{i=1}^{n}(L_{p,i}-\overline{L}_p)^2}{n(n-1)}}.
  \label{eq:niepewnosc}
\end{equation}
Dla danych demonstracyjnych $n=5$, $\overline{L}_p=\qty{70.0}{\decibel}$
i $u_A\approx\qty{0.07}{\decibel}$. Średnia arytmetyczna poziomów
opisuje tu estymację z powtórzeń; nie zastępuje uśredniania energetycznego
sygnałów zmiennych w czasie.

Uzupełnić budżet niepewności o składniki typu B wynikające z przyrządu,
kalibracji i metody. Przy niezależnych składnikach wyrażonych w jednostce
wyniku można zastosować
\begin{equation}
  u_c = \sqrt{u_A^2+\sum_j u_{B,j}^2},\qquad U=k u_c.
\end{equation}
Podać rozkłady, współczynniki wrażliwości i przyjęty współczynnik
rozszerzenia $k$. Nie przypisywać poziomu ufności bez uzasadnienia.
